% !TEX program = xelatex
% Residual UQ audit: Monte-Carlo error, FDR multiplicity, Wilcoxon discreteness.
% Tables generated by make_audit_tables.R.
\documentclass[12pt]{article}

\usepackage{fontspec}
\usepackage{xeCJK}
\setCJKmainfont[AutoFakeSlant=.2, AutoFakeBold=3]{AR PL KaitiM Big5}

\usepackage[margin=1in]{geometry}
\usepackage{amsmath, amssymb}
\usepackage{booktabs}
\usepackage{array}
\usepackage{hyperref}
\hypersetup{colorlinks=true, linkcolor=blue!45!black, urlcolor=blue!45!black}

\title{\textbf{殘餘不確定性稽核：MC error、FDR、Wilcoxon 離散}}
\author{清掉文獻檢視剩下的三條 (\#2--\#4)}
\date{2026 年 7 月}

\begin{document}
\maketitle

\noindent 承接 folder 08，本文把文獻檢視剩下的三條稽核項一次清掉：實測後驗
Monte-Carlo 誤差（先前只斷言）、對多重性做正式 FDR 校正（先前用嘴巴帶過）、
以及誠實面對 $n=10$ Wilcoxon 的 p 值離散。

\section*{\#2 後驗 Monte-Carlo 誤差：實測，不斷言}

folder 01 曾斷言「$\widehat P$ 的 MC 誤差 $\approx\sqrt{p(1-p)/M}$，遠小於評估
不確定性、可忽略」。用 batch-means（$5000$ 抽樣切 $10$ 批，批間 SD$/\sqrt{G}$
估全樣本 pooled 分數的 MC s.e.；兩鏈獨立故
$\mathrm{MC\_se}(\Delta)=\sqrt{\mathrm{se}_A^2+\mathrm{se}_B^2}$）實測後：

% Auto-generated by make_audit_tables.R
\begin{table}[htbp]\centering
\caption{Posterior Monte-Carlo error of the paired difference $\Delta$
(batch-means, $G{=}10$ batches of $500$ draws) against the site-bootstrap CI
half-width, ozone AR(2) vs.\ AR(1), $200/200$. The ratio is MC noise as a
fraction of evaluation uncertainty. Negligible it is not for plain CRPS.}
\label{tab:mc-error}\small
\begin{tabular}{@{}lcccc@{}}\toprule
Score & $\Delta$ & MC s.e.$(\Delta)$ & boot.\ CI half-width & ratio \\\midrule
Brier $q_{.95}$ & $-6.2e-05$ & $6.6e-05$ & $0.00035$ & $0.19$ \\
twCRPS (tail) & $-0.00015$ & $0.0012$ & $0.0045$ & $0.26$ \\
plain CRPS & $0.021$ & $0.013$ & $0.019$ & $0.70$ \\
\bottomrule\end{tabular}\end{table}


\noindent\textbf{斷言部分錯。}對 Brier 與 twCRPS，MC 噪音只佔評估不確定性的
$19\%$、$26\%$——次要，合併後 CI 僅寬 $2$--$3\%$，\textbf{null 不受影響}。但對
\textbf{plain CRPS，MC 噪音高達 $70\%$}，完全不可忽略：合併兩項
$\sqrt{0.0185^2+0.013^2}\approx0.023$ 已\emph{大於}效果量 $|\Delta|=0.021$。也就是
說，光是評估 $+$ MC 兩項就讓 CRPS 的「AR(2) 較差」不顯著——這是繼 twCRPS
（bulk）與 block bootstrap（不耐空間相依）之後，打死那格的\textbf{第三個角度}。

\section*{\#3 多重性：正式 FDR，不再用嘴巴帶過}

把全分析\emph{所有}配對比較的 p 值收成一個家族，套 Benjamini--Hochberg（BH）與
對相依穩健的 Benjamini--Yekutieli（BY）：

% Auto-generated by make_audit_tables.R
\begin{table}[htbp]\centering
\caption{Benjamini--Hochberg (BH) and Benjamini--Yekutieli (BY) FDR correction over the whole family of $m=117$ paired comparisons. 7 are nominally significant (raw $p<0.05$); \textbf{0 survive} either correction. The nominally significant cells are shown; all are plain-CRPS or the sparsest-level Brier, and none clears FDR.}
\label{tab:fdr}\small
\begin{tabular}{@{}llccc@{}}\toprule
Arm & Comparison & raw $p$ & $p_{\mathrm{BH}}$ & $p_{\mathrm{BY}}$ \\\midrule
ozone-200-200 & MRTS(10) vs none (brier) & $0.018$ & $0.458$ & $1.000$ \\
timeblock & set5 1-15 AR(2) - AR(1) (crps) & $0.021$ & $0.458$ & $1.000$ \\
ozone-200-200 & AR(2) vs AR(1) (crps) & $0.026$ & $0.458$ & $1.000$ \\
timeblock & set5 1-5 AR(2) - iid (crps) & $0.026$ & $0.458$ & $1.000$ \\
ozone-300-100 & AR(2) vs AR(1) (crps) & $0.027$ & $0.458$ & $1.000$ \\
timeblock & set5 1-5 AR(2) - AR(1) (crps) & $0.033$ & $0.458$ & $1.000$ \\
timeblock & set5 1-15 AR(2) - iid (crps) & $0.042$ & $0.458$ & $1.000$ \\
\bottomrule\end{tabular}\end{table}


\noindent$117$ 個比較中 $7$ 個名目顯著，\textbf{校正後 $0$ 個存活}（最小 raw
$p=0.018\to p_{\mathrm{BH}}=0.458$）。那 $7$ 格全是 plain-CRPS 或最稀疏門檻的
Brier——正是先前逐一質疑過的。手一揮的「多重性下不足採信」現在有形式化支撐：
\textbf{沒有任何一格通過 FDR $0.05$}。

\section*{\#4 Wilcoxon $n=10$：p 值是離散的，別假裝三位精度}

模擬與 time-block 用 $n=10$ 的 signed-rank，其 null 分布離散。可達的單尾 p 值在
$(0.005,0.10)$ 之間只有 $11$ 個：
\[
  0.0068,\,0.0098,\,0.0137,\,0.0186,\,\mathbf{0.0244},\,\mathbf{0.0322},\,
  \mathbf{0.0420},\,0.0527,\,0.0654,\,0.0801,\,0.0967,
\]
最小可達 $p=1/2^{10}=0.000977$。我先前（常態近似）報的 $0.026$、$0.033$、$0.042$
其實只能落在 $0.0244$、$0.0322$、$0.0420$——\textbf{報到小數三位是虛報精度}，且
這些「$<0.05$」的格正卡在門檻附近的網格點上，本就脆弱。（此點不改變任何結論——
經 \#3 校正後它們本來就都不顯著——但誠實報告應註明 $n=10$ 的離散性，或改用
exact test。）

\section*{總結：文獻檢視的每一條都已落地}

\begin{itemize}
  \item \textbf{\#1 站間相依}（folder 08）：null 全通過空間 block bootstrap；唯一顯著
        格瓦解。
  \item \textbf{\#2 MC 誤差}：Brier/twCRPS 次要（null 無恙）；CRPS 達 $70\%$，是打死
        那格的第三個角度。
  \item \textbf{\#3 多重性}：正式 FDR 下 $7$ 個名目顯著格\emph{全數}不存活。
  \item \textbf{\#4 Wilcoxon 離散}：報告精度已更正；不影響結論。
\end{itemize}

\noindent\textbf{淨效果}：論文核心結論（AR(2)/MRTS 在任何情境、任何尾部分數下都
不改變門檻超標的預測技巧）是 \textbf{null}，通過了站間相依、MC 誤差、多重性三重
加壓；而每一個曾經「顯著」的格，都在至少一個（多半是全部三個）角度下瓦解。整條
不確定性分析至此經得起教科書等級的檢視。

\medskip
\noindent\footnotesize 程式：\verb|ozone/US-all-auto/mc_error_check.R|、
\verb|simstudy/fdr_correction.R|。資料：\verb|mc_error_check.csv|、
\verb|fdr_correction.csv|。
\end{document}
